\documentclass[10pt,conference,compsocconf]{IEEEtran}

\usepackage{hyperref}
\usepackage{graphicx}
\usepackage{booktabs}

\begin{document}
\title{How Unified is Your Encoder?\\ Measuring Representational Convergence in Any-to-Any Vision Models}

\author{
  Arthur Margeat (330258), Adrien Clement (345535),
  Albert Fares (341018), Martina Gatti (341013)\\
  \textit{CS-503 Visual Intelligence --- Project Progress Report 1}
}

\maketitle

% ---------- Milestone Progress ----------
\section{Milestone Progress}

\subsection{Problem and milestone overview}
Any-to-any models such as 4M~\cite{4m} share Transformer encoder
weights across modalities, but route each modality through its own
learned embedding. Whether the encoder \emph{unifies} these inputs
into a shared representation, or behaves as an ensemble of disjoint
modality-specific sub-networks~\cite{zamir}, is open. RQ1 of our
proposal addresses this with a multi-metric protocol (CKA~\cite{cka},
PWCCA~\cite{pwcca}, k-NN overlap) layered over a mismatched-scene
null distribution, so that observed similarity is read against an
empirical noise trajectory rather than in absolute terms.

The first milestone was scoped around
\emph{building a strong, reusable benchmarking pipeline}, not around
breadth of modalities. Concretely, we restricted ourselves to a
Hypersim~\cite{hypersim} subset and a single modality pair
(RGB$\leftrightarrow$Depth), and instead invested effort into the
three-metric stack and the null-hypothesis machinery so that any
future pair plugs in without re-engineering. This deliberately trades
the breadth promised by the proposal (RGB, depth, normals,
segmentation) for confidence that the protocol itself is correct
before scaling. Extending to additional modalities and to
DIODE~\cite{diode} is the explicit goal of milestone~2; RQ2's
exploratory protocols (cross-modal transfer, residual analysis)
remain in the W5--W6 slot of the proposal timeline.

\subsection{Pipeline}
The benchmark proceeds in four steps. \\
(i)~Each modality of each scene is forwarded through 4M independently; forward hooks capture the output token matrix at each of the 12 encoder blocks, yielding per-layer activation matrices indexed by scene, modality, and layer. \\
(ii)~For every
layer and every modality pair, we compute CKA (global geometric
alignment), PWCCA (directional alignment of principal axes), and
k-NN overlap (local topological similarity); the three target
complementary facets of unification, so cross-metric agreement
sharpens any claim.\\ 
(iii)~Raw similarity is confounded by shared
positional encodings, tokenizer priors, and inter-scene statistical
We construct a layer-wise empirical null by recomputing each metric on $1{,}000$ mismatched cross-scene modality pairs, yielding a noise trajectory across depth. To prevent the null from being inflated by scenes that share the same semantic content, left and right scenes are constrained to have \emph{different scene types} (e.g.\ bathroom vs.\ office) drawn from the Hypersim trajectory metadata; this ensures the null reflects genuine cross-content misalignment rather than intra-category similarity. \\
(iv)~Matched values are tested
against this null per layer: we report p-values, null mean and
standard deviation, and the matched$-$null gap.\\
The four steps are
decoupled ---the comparison step re-runs without re-extracting
activations--- and all artefacts are logged to W\&B.

\subsection{Experiment}
We ran the pipeline end-to-end on a Hypersim subset, $4325$ aligned
RGB$\leftrightarrow$Depth samples, with $1000$ mismatched null draws
across the $12$ encoder blocks of 4M-7B-CC12M. All three metrics
produced enriched outputs and the corresponding plots.

\begin{figure}[t]
  \centering
  \includegraphics[width=\columnwidth]{figures/metric_vs_layer.pdf}
  \caption{Layer-wise convergence on Hypersim
  (RGB$\leftrightarrow$Depth, $n{=}4325$, $1000$ null draws).}
  \label{fig:metric-vs-layer}
\end{figure}

\begin{figure}[t]
  \centering
  \includegraphics[width=0.65\columnwidth]{figures/pvalue_vs_layer.pdf}
  \caption{Per-layer p-values for the three metrics. k-NN overlap drops well below threshold from layer~$7$
  onward; CKA hovers around it; PWCCA never reaches significance.}
  \label{fig:pvalue-vs-layer}
\end{figure}

\subsection{Results}
Two of the three metrics agree on a clear depth-wise pattern
(Fig.~\ref{fig:metric-vs-layer}, Fig.~\ref{fig:pvalue-vs-layer}).
\textbf{CKA} rises monotonically from $0.10$ at layer~$0$ to $0.19$ at
layer~$10$ ($z{=}2.60$, $p{=}0.028$), then dips to $0.17$ at the final
block. \textbf{k-NN overlap} shows the strongest signal: $0.09$ at
layer~$0$ grows to $0.35$ at layer~$9$ ($z{=}3.62$,
$p{\approx}10^{-3}$) and remains highly significant through
layer~$11$. Both shapes match the ``persistent partial convergence''
outcome anticipated by Fig.~2 of our proposal: significance grows with
depth and the mild retreat at the last block is consistent with
task/decoder specialisation. Crucially, CKA peaks at $0.19$, far from
$1$, supporting the proposal's working definition of \emph{partial}
unification rather than a fully shared trunk: the encoder pulls RGB
and depth toward a common subspace without collapsing them onto it.

\textbf{PWCCA} disagrees. Observed values ($0.25$--$0.56$) lie below
the null mean ($\approx 0.70$) at every layer; all p-values exceed
$0.8$, z-scores are negative. We interpret this as a metric pathology in our setup; one plausible explanation is that PWCCA's CCA backbone recovers spurious shared subspaces driven by tokenizer and positional artefacts that are shared across scenes regardless of content, inflating the null relative to truly matched pairs. Confirming this diagnosis is a milestone-2 priority.

\subsection{Deviations from the proposal}
Two deliberate scope changes versus the proposal: (i) the milestone-1
pipeline run uses one modality pair instead of all four, on a Hypersim
subset; the breadth is deferred to milestone~2 once the protocol is
trusted. (ii) PWCCA is uninformative under our current null
construction, so its weight in the final claim will likely shift onto
CKA and k-NN overlap unless we recover it via a stricter null. The
research question, datasets, and overall plan are unchanged.

% ---------- Discussion & Next Steps ----------
\section{Discussion \& Next Steps}
\label{sec:next}

\paragraph{Reading}
Across two of three metrics the encoder shows depth-correlated,
statistically credible alignment between RGB and depth that peaks at
layers~$9$--$10$ and retreats at layer~$11$. In the language of the
proposal's working definition, this is evidence for \emph{partial}
unification with a specialised tail---neither the ``red'' fragmented
outcome of Fig.~2 nor a fully unified encoder.

\paragraph{What worked}
The pipeline is functional, optimised end-to-end, and lightweight
enough to be fully replicable: a complete run with $1000$ null draws
finishes in minutes on a single GPU, and re-running only the
significance step takes seconds.

\paragraph{What did not work, and why}
PWCCA never reaches significance, and observed values sit
\emph{below} the null mean at every layer. We do not yet have a
confident explanation; diagnosing whether this stems from the metric,
the null construction, or the data is on the milestone-2 list.

\paragraph{Risks}
(i) \emph{Hypersim is synthetic}; convergence on real images may
differ, motivating the DIODE replication. (ii) \emph{One pair only}
cannot distinguish ``RGB and depth converge'' from ``everything
converges''; resolved by adding modalities. (iii) \emph{Scale}: the
current run uses a small Hypersim subset, so signal stability and
generalisation must be verified on a substantially larger sample
before any conclusion is treated as final.


\paragraph{Action items until the next report}
\begin{enumerate}\itemsep0pt
  \item Scale the Hypersim run up to a much larger subset to verify
        the depth trend and tighten the null estimates.
  \item Replicate on DIODE (real-world RGB$+$depth) using the
        existing adapter, isolating the synthetic-data risk.
  \item Add surface normals and segmentation as additional
        modalities, yielding more pairs and a transitivity check on
        convergence.
  \item Diagnose PWCCA: identify why it underperforms the null and
        decide whether to keep, replace, or drop it.
  \item Begin RQ2 explorations (cross-modal $k$-NN, residual
        correlation) on the W5--W6 slot of the proposal timeline.
\end{enumerate}

% ---------- Author Contribution Statement ----------
\section{Author Contribution Statement}
\textit{Arthur} led the extraction stage and the 4M model wrapper.
\textit{Albert} and \textit{Martina} implemented the three metrics
(CKA, PWCCA, k-NN overlap) and the null-hypothesis stage. \textit{Adrien} contributed to result
analysis and the write-up. All authors participated in shaping the
research question.


\begin{thebibliography}{99}\small
\bibitem{4m} D.~Mizrahi \emph{et al.}, ``4M: Massively multimodal masked modeling,'' \emph{NeurIPS}, 2023.
\bibitem{zamir} A.~Zamir, EPFL Actu interview, 2025.
\bibitem{cka} S.~Kornblith \emph{et al.}, ``Similarity of neural network representations revisited,'' \emph{ICML}, 2019.
\bibitem{pwcca} A.~Morcos, M.~Raghu, S.~Bengio, ``Insights on representational similarity \ldots'' \emph{NeurIPS}, 2018.
\bibitem{hypersim} M.~Roberts, N.~Paczan, ``Hypersim,'' \emph{ICCV}, 2021.
\bibitem{diode} I.~Vasiljevic \emph{et al.}, ``DIODE,'' arXiv:1908.00463, 2019.
\end{thebibliography}

\end{document}
